\section*{Chapitre \ref{ch:iv1:games} --- Jeux}
\markboth{\protect\omTocMark{}{Réponses~: Jeux}}{\protect\omTocMark{}{Réponses~: Jeux}}

\begin{solution}{iq:iv1:games:1}
Jouer le premier coup et prendre 6, ce qui laisse 91. Si prendre la dernière pierre fait perdre, jouer le premier coup et prendre 5, ce qui laisse 92.

Les multiples de 7 sont des positions perdantes pour le joueur qui doit jouer\index{positions perdantes}. Depuis un multiple de 7, tout coup mène à un non-multiple~; depuis un non-multiple, on peut toujours revenir à un multiple. Ainsi, quoi que prenne l'adversaire, on complète à 7. Comme $97 = 13 \cdot 7 + 6$, on prend 6.

Si la dernière pierre fait perdre, on veut que l'adversaire se retrouve avec exactement 1 pierre. Les positions perdantes sont donc celles situées juste au-dessus d'un multiple de 7, et il faut laisser 92.

\iqwrong{Affirmer que le total perdant est le même dans les deux versions. La cible se décale d'un cran, car la dernière pierre est désormais une pierre à éviter.}
\end{solution}

\begin{solution}{iq:iv1:games:2}
Le premier joueur gagne en prenant 3 pierres dans le tas le plus grand, ce qui ramène les deux tas à 5. Si les tas sont égaux au départ, c'est le second joueur qui gagne.

Des tas égaux sont perdants pour le joueur qui doit jouer\index{stratégie miroir}. Quoi que ce joueur prenne dans un tas, l'adversaire prend autant dans l'autre, de sorte que les tas restent égaux et que l'adversaire obtient la dernière pierre. Des tas inégaux sont gagnants, car un seul coup permet de les égaliser. Une table établie par recherche exhaustive, jusqu'à des tas de 9 pierres, confirme la règle.

\iqwrong{Penser que le plus grand tas est un avantage en soi, et chercher à le conserver. Seul compte l'écart, et le gagnant est celui qui laisse un écart de 0.}
\end{solution}

\begin{solution}{iq:iv1:games:3}
Chaque joueur montre face avec la probabilité $3/8$, et le jeu vaut $-\$1/8$ par manche pour le joueur considéré.

Un mélange doit rendre l'autre joueur indifférent\index{stratégie mixte}, sinon celui-ci l'exploiterait. Si le premier joueur montre face avec la probabilité $q$, les deux choix du second joueur rapportent
\[
3q - 2(1 - q) \quad\text{and}\quad -2q + (1 - q),
\]
et ces deux expressions sont égales pour $q = 3/8$, ce qui donne une valeur de $-1/8$. Le mélange $p$ du second joueur résout les mêmes équations, et il vaut aussi $3/8$.

Une seconde méthode consiste à additionner les quatre issues sous ces mélanges~: $\tfrac{9}{64} \cdot 3 + \tfrac{25}{64} \cdot 1 - \tfrac{30}{64} \cdot 2 = -\tfrac18$.

Les quatre gains s'additionnent à zéro, et le jeu semble donc équitable, mais il ne l'est pas. Le gros gain est celui que l'on encaisse rarement~: face contre face n'arrive que dans $9/64$ des cas, car le premier joueur l'évite.

\iqwrong{Dire que chacun devrait montrer face plus souvent, puisque le gros gain vient de face. Si l'on montre face davantage, l'adversaire montre pile, et l'on paie 2~\$ plus souvent. Les gains de chacun fixent le mélange de l'adversaire, et non le sien.}
\end{solution}

\begin{solution}{iq:iv1:games:4}
Des joueurs parfaitement rationnels écrivent 0, ou 1, puisque avec des nombres entiers les deux sont stables. Dans une salle réelle, on écrit un nombre supérieur à cela, car les autres ne descendent pas jusqu'au bout.

Les deux tiers de la moyenne ne peuvent jamais dépasser $66.7$, donc tout nombre supérieur à 67 ne fait pas mieux que 67, quoi que les autres écrivent\index{dominance itérée}. Si tout le monde le sait, la moyenne est au plus égale à 67, donc aucun nombre au-dessus de 45 ne vaut la peine d'être écrit. En répétant le raisonnement, on ramène le plafond à 1. Si tous écrivent 0, la cible est 0. Si tous écrivent 1, la cible est $\tfrac23$. Dans une salle de cinq personnes ou plus, la cible reste plus proche de 1 que de 0 même après qu'une personne est passée à 0, donc personne n'a intérêt à descendre.

Dans une salle réelle, la question est de savoir combien d'étapes les autres franchissent. S'ils choisissent au hasard, en moyenne 50, la meilleure réponse est 33. Si l'on pense qu'ils sont tous à 33, la réponse est 22, puis 15. Il faut donc choisir un nombre bas, situé d'une étape au-delà de ce que font les autres, et non la réponse théorique.

\iqwrong{Écrire 0 et revendiquer la victoire. C'est l'équilibre si tout le monde est rationnel et le sait, mais la question demande de gagner. Face à de vraies personnes, trop raisonner fait perdre autant que trop peu raisonner.}
\iqwhy{C'est un modèle de marchés où la valeur dépend de ce que les autres pensent que les autres pensent. La réponse rationnelle et la réponse profitable peuvent différer, et avoir une étape d'avance l'emporte sur avoir raison.}
\end{solution}

\begin{solution}{iq:iv1:games:5}
Jouer la pierre dans $1/2$ des cas, le papier dans $1/6$ et les ciseaux dans $1/3$. La pierre est jouée le plus souvent, alors qu'une victoire avec la pierre rapporte le moins.

Dans un jeu symétrique à somme nulle comme celui-ci, la valeur est 0 et chaque choix doit rapporter 0 face au mélange. Soit le mélange $(r, p, s)$\index{stratégie mixte}~:
\[
\text{pierre~: } s - 2p = 0, \quad \text{papier~: } 2r - 3s = 0, \quad \text{ciseaux~: } 3p - r = 0.
\]
Donc $r = 3p$ et $s = 2p$, et leur somme vaut $6p = 1$.

Cela paraît étrange, car on s'attendrait à ce que le choix le plus rémunérateur soit le plus fréquent. Mais les fréquences ne sont pas fixées par ce qu'un choix rapporte à celui qui le joue. Elles sont fixées par ce qu'il fait à l'autre joueur. La pierre est fréquente parce qu'elle doit dissuader le choix qui rapporte le plus~: si la pierre était rare, les ciseaux rapporteraient leurs 3~\$ sans aucun risque. Un programme linéaire sur la matrice $3 \times 3$ donne le même mélange.

\iqwrong{Mélanger en proportion des gains, $1:2:3$. Face à ce mélange, la pierre rapporte $\tfrac36 \cdot 1 - \tfrac26 \cdot 2 = -\tfrac16$, le papier rapporte $\tfrac16 \cdot 2 - \tfrac36 \cdot 3 = -\tfrac76$, et les ciseaux rapportent $\tfrac26 \cdot 3 - \tfrac16 \cdot 1 = \tfrac56$, donc l'adversaire joue systématiquement les ciseaux.}
\end{solution}

\begin{solution}{iq:iv1:games:6}
Jouer en premier. La somme de Nim vaut $3 \oplus 5 \oplus 7 = 1$, qui n'est pas nulle. Il existe trois coups gagnants~: vers $(2, 5, 7)$, $(3, 4, 7)$ et $(3, 5, 6)$.

La règle\index{Nim}\index{somme de Nim}~: une position est perdante pour le joueur qui doit jouer exactement lorsque le ou exclusif bit à bit des tas vaut 0. Voici la démonstration, en trois étapes.

D'abord, la position vide a une somme de Nim nulle, et elle est perdante.

Ensuite, tout coup depuis une somme de Nim nulle donne une somme non nulle, car modifier un tas change ses bits, les autres restant inchangés.

Enfin, depuis une somme de Nim non nulle $s$, il existe un coup vers 0. Prendre un tas dont le bit de poids fort de $s$ vaut 1, et le remplacer par le tas $\oplus\ s$. Cette nouvelle valeur est plus petite, puisque le bit de poids fort de $s$ passe de 1 à 0, et le total devient 0.

Ici $s = 1$, dont le bit de poids fort est le bit 0. Les trois tas sont impairs, donc chacun peut servir~: $3 \to 2$, $5 \to 4$, $7 \to 6$. Une recherche exhaustive sur tous les tas jusqu'à 5, ainsi que sur cette position, confirme la règle et les trois coups.

\iqwrong{Chercher à raisonner sur la parité du nombre de pierres, ou à garder les tas équilibrés en taille. Ce sont les chiffres binaires qui décident, et des tas de $1, 2, 3$ pierres sont perdants pour le joueur qui doit jouer, alors qu'ils paraissent déséquilibrés.}
\end{solution}

\begin{solution}{iq:iv1:games:7}
Les valeurs de Grundy pour 0 à 8 pierres sont $0, 1, 0, 1, 2, 3, 2, 0, 1$, de période 7. Le premier joueur gagne avec des tas de 10 et 6, car $g(10) \oplus g(6) = 1 \oplus 2 = 3$. Les coups gagnants consistent à prendre 4 dans le tas de 10, ce qui laisse 6 et 6, ou à prendre 3 dans le tas de 6, ce qui laisse 10 et 3.

La valeur de Grundy d'une position est le plus petit entier naturel qui ne figure pas parmi les valeurs de ses coups\index{valeur de Grundy}. Une position de valeur $v$ se comporte, dans une somme, comme un tas de Nim de taille $v$\index{théorème de Sprague--Grundy}. Ainsi, pour $n = 4$, les coups mènent à 3, 1 et 0, de valeurs $1, 1, 0$, donc $g(4) = 2$. Pour $n = 5$, les coups mènent à 4, 2, 1, de valeurs $2, 0, 1$, donc $g(5) = 3$.

Une somme est perdante lorsque la somme de Nim des valeurs est nulle. Il faut un coup dans un tas qui rende sa valeur égale à celle de l'autre. Depuis 10, une valeur de 2 est atteignable en prenant 4. Depuis 6, une valeur de 1 est atteignable en prenant 3, car $g(3) = 1 = g(10)$. Une table exhaustive, établie sur les couples de tas, trouve ces deux coups et aucun autre.

\iqwrong{Appliquer le Nim simple aux tailles des tas, $10 \oplus 6 = 12$. Les tailles des tas ne conviennent que si l'on peut retirer n'importe quel nombre de pierres. Ici, il faut utiliser les valeurs de Grundy.}
\end{solution}

\begin{solution}{iq:iv1:games:8}
A propose 98 pour A, 0 pour B, 1 pour C, 0 pour D et 1 pour E. Si une proposition doit recueillir plus de la moitié des voix, A propose 97 pour A, 1 pour C et 2 pour l'un de D ou E.

Raisonner à rebours\index{induction à rebours}\index{jeu des pirates}, en supposant qu'au moins la moitié des voix soit nécessaire.

Avec D et E, D vote pour son propre plan, qui représente la moitié des voix~: $(100, 0)$.

Avec C, D et E~: C a besoin d'une voix supplémentaire. E obtient 0 sinon, donc E se contente de 1~: $(99, 0, 1)$.

Avec B, C, D et E~: B a besoin d'une voix de plus. D obtient 0 sinon~: $(99, 0, 1, 0)$.

Avec les cinq~: A a besoin de deux voix de plus, provenant de ceux qui obtiendraient 0 sous la proposition de B, à savoir C et E~: $(98, 0, 1, 0, 1)$.

Lorsqu'il faut plus de la moitié des voix, D est condamné dès lors que seuls D et E restent, puisque E peut rejeter et tout garder. D soutiendra donc C sans contrepartie, et C garde 100 avec la voix de D. B a besoin de trois voix sur les quatre, et D et E reçoivent chacun 1, soit $(98, 0, 1, 1)$. A a besoin de trois voix sur cinq, et les moins chères sont C, pour 1, et l'un de D ou E, pour 2.

\iqwrong{Donner à chaque pirate une part égale pour obtenir des voix. Un votant se compare au tour suivant, et non aux autres, et une voix coûte une pièce de plus que ce que ce votant obtiendrait au tour suivant.}
\iqwhy{Dans tout vote ou toute négociation, le prix d'un votant dépend de ce qu'il obtient si la proposition en cours échoue. C'est la logique d'une position de négociation.}
\end{solution}

\begin{solution}{iq:iv1:games:9}
Avec trois tours, A obtient $0.91$. Sans limite, A obtient $1/(1 + 0.9) = 10/19 \approx 0.526$.

Pour trois tours, on raisonne à rebours\index{induction à rebours}, dans la monnaie de ce tour. Au tour 3, A ne propose rien à B et garde 1, car B n'obtient rien en refusant. Au tour 2, B doit offrir à A ce que le tour 3 vaut pour A aujourd'hui, soit $0.9$, et B garde $0.1$. Au tour 1, A doit offrir à B $0.9 \cdot 0.1 = 0.09$, et A garde $0.91$.

Sans limite, soit $x$ ce que obtient le proposant\index{négociation par offres alternées}. Le jeu reste le même après un refus, donc le répondant doit recevoir $\delta x$, et le proposant garde $1 - \delta x$. Alors $x = 1 - \delta x$, ce qui donne $x = \frac{1}{1 + \delta}$. Reprendre l'induction à rebours sur un grand nombre de tours converge vers ce résultat.

Les limites ont un sens. Pour $\delta = 0$, le proposant prend tout, et lorsque $\delta \to 1$, chacun obtient la moitié. Proposer en premier vaut davantage lorsque le délai coûte cher.

\iqwrong{Utiliser le résultat du dernier à proposer pour le jeu sans limite. La large part de $0.91$ vient de ce que A a la dernière offre, et sans dernier tour, cet avantage disparaît.}
\end{solution}

\begin{solution}{iq:iv1:games:10}
Le tireur vise à gauche avec la probabilité $8/15$, le gardien plonge à gauche avec la probabilité $3/5$, et la probabilité d'un but vaut $0.74$. Après l'amélioration, le tireur vise à droite moins souvent~: $7/16$ au lieu de $7/15$.

Le mélange du tireur rend le gardien indifférent\index{stratégie mixte}. Avec $p$ la probabilité de viser à gauche, la probabilité de but vaut
\[
0.6p + 0.9(1 - p) \text{ si le gardien plonge à gauche, et } 0.95p + 0.5(1 - p) \text{ si à droite.}
\]
Ces deux expressions coïncident pour $p = 8/15$, avec une valeur de $0.74$. La probabilité $q$ que le gardien plonge à gauche vérifie $0.6q + 0.95(1 - q) = 0.9q + 0.5(1 - q)$, donc $q = 3/5$.

Lorsque le $0.9$ passe à $0.95$, la première expression change, et $p = 9/16$. Viser à droite recule donc de $7/15 \approx 0.467$ à $7/16 \approx 0.438$, et la valeur monte à environ $0.753$.

La surprise est juste. Le mélange du tireur doit maintenir le gardien indifférent, et lorsque viser à droite est devenu plus dangereux, le gardien plonge moins souvent à gauche. Ce déplacement du gardien rend viser à gauche plus avantageux. L'option améliorée est utilisée moins souvent, et le tireur gagne tout de même. Un programme linéaire donne ces deux résultats.

\iqwrong{Dire qu'il faut utiliser davantage sa compétence améliorée. Dans un jeu avec un adversaire qui réagit, un changement de ses gains déplace le mélange de l'adversaire, et le sien est fixé par ce qui rend l'adversaire indifférent.}
\iqwhy{Un signal amélioré ou une plateforme de négociation améliorée ne signifie pas toujours qu'il faut l'utiliser davantage. Si les autres joueurs réagissent, le résultat relève de l'équilibre, et non de votre propre tableau de gains.}
\end{solution}

\begin{solution}{iq:iv1:games:11}
Dans une enchère au second prix, enchérissez votre valeur. Dans une enchère au premier prix, enchérissez $\frac{n - 1}{n}v$. Dans les deux cas, le vendeur espère $\frac{n - 1}{n + 1}$.

Au second prix~: votre enchère décide seulement si vous gagnez, jamais ce que vous payez\index{enchère au second prix}. Si la plus haute enchère des autres est inférieure à votre valeur, vous voulez gagner, et si elle est supérieure, vous ne le voulez pas. Enchérir votre valeur fait les deux, donc elle est dominante.

Au premier prix~: supposons que les autres enchérissent $\beta v$\index{enchère au premier prix}. Enchérir $b$ gagne lorsque les $n - 1$ autres ont des valeurs inférieures à $b/\beta$, donc le gain vaut $(v - b)(b/\beta)^{n - 1}$. En annulant la dérivée, on obtient $b = \frac{n - 1}{n}v$, pour tout $\beta$, et cela est donc cohérent. Un balayage des enchères pour $v = 0.9$ confirme le maximum pour $n = 2, 3, 5$.

Recette~: l'enchère au second prix rapporte la deuxième plus haute valeur, dont la moyenne vaut $\frac{n - 1}{n + 1}$. L'enchère au premier prix rapporte $\frac{n - 1}{n}$ fois la plus haute valeur, dont la moyenne vaut $\frac{n}{n + 1}$, ce qui donne le même résultat. C'est l'équivalence des recettes\index{équivalence des recettes}, et il ne s'agit que d'une égalité des moyennes. La recette de l'enchère au premier prix a une variance plus faible dans la simulation.

\iqwrong{Supposer qu'une enchère au premier prix rapporte davantage parce que le gagnant paie sa propre enchère. Les enchérisseurs minorent leurs enchères pour en tenir compte, et dans ce cadre la minoration compense exactement la différence.}
\end{solution}

\begin{solution}{iq:iv1:games:12}
Chacun le fait avec la probabilité $3/4$ pour $n = 2$ et $1/2$ pour $n = 3$. Personne ne le fait avec une probabilité de $1/16$ et $1/8$, et cette probabilité croît vers $1/4$ lorsque $n$ augmente, tandis que la probabilité propre à chacun tend vers 0.

Une stratégie mixte\index{stratégie mixte} n'est utilisée que si les deux choix rapportent la même chose. Le faire rapporte 3 avec certitude. Ne pas le faire rapporte 4 si au moins un des $n - 1$ autres le fait. Avec la probabilité $p$ pour chacun, la condition\index{dilemme du volontaire} est
\[
3 = 4\bigl(1 - (1 - p)^{n - 1}\bigr), \qquad\text{soit}\qquad (1 - p)^{n - 1} = \tfrac14.
\]
Donc $p = 3/4$ pour $n = 2$ et $p = 1/2$ pour $n = 3$. La probabilité que personne ne le fasse vaut $(1 - p)^n = (\tfrac14)^{n/(n - 1)}$, ce qui donne $1/16$, $1/8$, et tend vers $1/4$.

Le gain de chacun vaut 3 quel que soit $n$, car le faire est toujours l'un de ses meilleurs choix. Un groupe plus grand n'aide donc personne ici.

\iqwrong{Dire qu'avec davantage de personnes, la tâche a plus de chances d'être accomplie. Chacun compte davantage sur les autres à mesure que le groupe grandit, et cet effet l'emporte sur les bras supplémentaires. Qu'une seule personne le fasse constitue aussi un équilibre, mais seulement si le groupe peut convenir de qui s'en charge.}
\iqwhy{Un risque que n'importe qui dans un desk pourrait signaler, à faible coût privé, a cette forme. Un groupe plus grand n'est pas plus sûr, sauf si quelqu'un est désigné comme responsable.}
\end{solution}

\begin{solution}{iq:iv1:games:13}
Le premier joueur gagne, par un argument de vol de stratégie\index{vol de stratégie}. Sur une barre carrée, le coup gagnant explicite mange le carré en diagonale voisin de celui du coin inférieur gauche, ce qui laisse deux bras égaux~; ensuite, le premier joueur reproduit les coups de l'adversaire en miroir.

Supposons que le premier joueur ne mange que le carré en haut à droite. Si cela laisse à l'adversaire une position perdante, le premier joueur a gagné. Sinon, le second joueur dispose d'une réponse gagnante, qui mange un certain carré $(i, j)$. Mais manger $(i, j)$ mange aussi le carré en haut à droite, car celui-ci se situe au-dessus et à droite de chaque carré. Le premier joueur aurait donc pu jouer $(i, j)$ directement et atteindre la même position. Dans les deux cas, il existe un premier coup gagnant. La preuve n'est pas constructive, et elle ne dit rien sur le coup lui-même.

Carré~: manger le carré voisin en diagonale du coin. Il reste une rangée du bas et une colonne de gauche de même longueur. Quoi que l'adversaire mange sur un bras, manger la même chose sur l'autre. L'adversaire se retrouve obligé de manger le coin.

Un arbre de jeu par force brute pour toutes les barres jusqu'à 5 par 5 confirme que le premier joueur gagne toutes les barres sauf la $1 \times 1$.

\iqwrong{Chercher un motif pour le coup gagnant en général. Aucune formule simple n'est nécessaire ici, et la question porte sur l'existence. Il faut le dire et donner l'argument.}
\end{solution}

\begin{solution}{iq:iv1:games:14}
Avec 201 pirates, le proposant survit sans pièce, et avec 202 de même. Avec 203, le proposant meurt, et avec 204, il survit. Au-delà, le proposant ne survit que pour $200 + 2^k$ pirates~: 208, 216, 232, 264, 328, et ainsi de suite. Avec 300 pirates, 36 sont jetés par-dessus bord.

Les pièces permettent d'acheter au plus 100 voix\index{problème des pirates}. Avec 200 pirates, le proposant a besoin de 100 voix, en paie 99 et garde 1. Avec 201, il faut 101 voix~: 100 achetées, plus la sienne, ce qui ne laisse aucune pièce. Avec 202, il faut encore 101 voix, et le même plan fonctionne.

Avec 203, il faut 102 voix, et seulement 101 voix sont disponibles. Le proposant meurt donc, quel que soit le plan. Cela crée une nouvelle source de voix~: un proposant qui mourrait au tour suivant votera pour n'importe quoi, afin de survivre.

Avec 204, il faut 102 voix~: 100 achetées, celle du proposant, et celle du suivant, qui mourrait sinon.

En général, soit $m$ la dernière taille à laquelle le proposant survit. Pour une taille $n = m + j$, le proposant dispose de $j$ voix qui assurent sa survie, la sienne comprise, plus 100 voix achetées. Il lui faut $\lceil n/2 \rceil$, ce qui donne $j \ge m - 200$. Ainsi, à partir de $m = 202$, on obtient 204, puis 208, 216 et ainsi de suite, et chaque survivant se situe à une distance du seuil 200 deux fois plus grande.

Pour 300, le survivant immédiatement inférieur est 264, donc les 36 pirates du haut meurent les uns après les autres. Avec 299, il faut 150 voix, et il y en a $35 + 100$ disponibles~; le manque subsiste jusqu'à 264.

\iqwrong{Prolonger simplement l'allocation des pièces au-delà de 200, ou penser que tout le monde meurt une fois les pièces épuisées. La peur d'être le suivant crée des voix gratuites, et c'est pourquoi certaines des tailles plus grandes fonctionnent.}
\end{solution}

\begin{solution}{iq:iv1:games:15}
X mise toujours l'as, passe avec le roi, et bluffe la dame avec la probabilité $1/3$. Y suit avec l'as, suit avec le roi avec la probabilité $1/3$, et se couche avec la dame. Le jeu vaut $1/18$ de dollar pour X, soit environ $0.056$.

Certains choix sont dominés\index{fréquence de bluff}. Avec l'as, X doit miser, et Y doit suivre. Avec la dame, Y doit se coucher. Avec le roi, miser est mauvais pour X~: la dame se couche, ce que le passage rapporterait de toute façon, et l'as suit, ce qui coûte davantage. Il ne reste donc que deux nombres.

La fréquence de bluff $b$ rend Y indifférent avec un roi. Lorsque Y détient le roi et fait face à une mise, X a l'as (1 façon) ou une dame qui bluffe ($b$ façons). Suivre rapporte $\frac{-2 + 2b}{1 + b}$, et se coucher rapporte $-1$. Ces deux valeurs coïncident pour $b = 1/3$.

La fréquence de suivi $c$ rend X indifférent avec une dame. Lorsque X détient la dame, passer perd 1. Miser contre le roi rapporte $1 - c$ lorsque Y se couche et $-2c$ lorsque Y suit, et miser contre l'as perd 2. Donc $-1 = \tfrac12\bigl((1 - c) - 2c\bigr) + \tfrac12(-2)$, ce qui donne $c = 1/3$.

La valeur provient d'un programme linéaire sur le tableau $8 \times 8$ des stratégies pures, soit $1/18$, et une simulation des deux mélanges confirme ce résultat. X l'obtient parce que lui seul peut décider si davantage d'argent est engagé, de sorte que Y doit payer pour le droit de repérer un bluff. Le faible avantage vient de ce choix, et non de meilleures cartes, puisque la donne est équitable pour les deux joueurs.

\iqwrong{Bluffer une fois sur deux, ou jamais. Si vous ne bluffez jamais, le roi se couche toujours, et votre as rapporte moins qu'il ne devrait. Si vous bluffez trop, le roi suit toujours. La bonne fréquence rend indifférent celui qui suit. Suivre coûte un dollar pour gagner un pot de trois, donc un bluff pour trois mises de valeur.}
\end{solution}

\begin{solution}{iq:iv1:games:16}
La meilleure réserve est $1/2$, avec une recette de $5/12$, contre $1/3$ sans réserve. Elle ne dépend pas du nombre d'enchérisseurs. Une réserve réduit le surplus total, de $2/3$ à $7/12$.

Enchérir sa valeur reste dominant. Si les deux valeurs sont supérieures à $r$, le prix est la plus basse. Si une seule l'est, le prix vaut $r$. Si aucune ne l'est, il n'y a pas de vente\index{prix de réserve}. Donc
\[
R(r) = \int_r^1 x \cdot 2(1 - x)\,dx + r \cdot 2r(1 - r) = \frac13 + r^2 - \frac43 r^3.
\]
En posant $R'(r) = 2r - 4r^2 = 0$, on obtient $r = \tfrac12$, et $R(\tfrac12) = 5/12$. Une simulation le confirme.

Pourquoi ne dépend-elle pas de $n$~? La réserve est une règle pour une seule vente. Elle ne joue que lorsqu'un enchérisseur est au-dessus et tous les autres en dessous, et l'on fixe alors en pratique un prix pour un acheteur unique dont la valeur est uniforme. Un acheteur dont la valeur est supérieure à $r$ paie $r$ avec la probabilité $1 - r$, et le produit est maximal en $\tfrac12$, quels que soient les autres participants. Avec 4 enchérisseurs, un balayage des réserves donne le même résultat.

Pour le surplus, le bien revient au plus offrant, sauf si les deux valeurs sont inférieures à $\tfrac12$, ce qui arrive avec la probabilité $\tfrac14$. Le surplus vaut alors $\int_{1/2}^1 x \cdot 2x\,dx = 7/12$, au lieu de $E[\max] = 2/3$.

\iqwrong{Dire qu'une réserve ne peut que nuire, parce qu'elle risque de ne pas vendre. Elle coûte $\tfrac14$ des ventes, mais elle relève le prix payé dans les autres cas, et le vendeur y gagne.}
\end{solution}

\begin{solution}{iq:iv1:games:17}
Le premier joueur gagne, par un seul coup~: prendre 2 pièces de la marche~1, ce qui laisse $(1, 1, 4, 1, 5)$. Seules comptent les marches impaires, et leur somme de Nim vaut $3 \oplus 4 \oplus 5 = 2$.

On traite les marches impaires comme des tas de Nim\index{Nim en escalier}. Déplacer des pièces d'une marche impaire vers la marche située en dessous est un coup de Nim sur ce tas, puisque les pièces arrivent sur une marche paire, qui ne compte pas. Déplacer des pièces depuis une marche paire les fait arriver sur une marche impaire, ce qui agrandit un tas. Mais l'adversaire peut toujours annuler cet effet en déplaçant les mêmes pièces plus bas, depuis cette marche impaire vers la marche paire située en dessous (depuis la marche~1, elles quittent l'escalier, ce qui revient au même). Un coup depuis une marche paire ne rapporte donc rien. La partie doit finir, car les pièces ne descendent que.

Le joueur qui se trouve face à une somme de Nim nulle sur les marches impaires perd. Ici, il faut la ramener à 0. Avec $s = 2$, le tas dont le bit~1 vaut~1 est 3, qui devient $3 \oplus 2 = 1$. Il faut donc prendre 2 pièces de la marche~1. Les coups depuis les marches paires n'aident pas~: ils donnent des sommes égales à 5 ou à 3.

Une recherche exhaustive de tous les coups depuis cette position trouve ce seul coup, et la règle coïncide avec la recherche exhaustive sur tous les petits escaliers.

\iqwrong{Compter les cinq tas dans la somme de Nim. Les marches paires peuvent toujours être annulées, donc ce ne sont pas de vrais choix.}
\end{solution}

\begin{solution}{iq:iv1:games:18}
Les positions perdantes sont $(0, 0)$, $(1, 2)$, $(3, 5)$, $(4, 7)$, $(6, 10)$, $(8, 13)$ et $(9, 15)$. La position $(9, 15)$ perd pour le joueur qui doit jouer, et $(10, 15)$ gagne, en prenant 1 au premier tas. La forme close est $\bigl(\lfloor k\varphi \rfloor, \lfloor k\varphi^2 \rfloor\bigr)$, avec $\varphi = (1 + \sqrt{5})/2$\index{jeu de Wythoff}.

On construit la liste à la main. Une position perdante ne peut pas partager la taille d'un tas avec une autre position perdante, puisqu'un coup sur un tas les relierait. Elle ne peut pas non plus partager une différence avec une autre, car un coup sur les deux tas conserve la différence. Ainsi, à l'étape $k$, on prend le plus petit entier strictement positif pas encore utilisé comme petit tas, et l'on ajoute $k$ pour obtenir le grand. Toute autre position peut atteindre une position perdante en un coup, comme le confirme une table de recherche exhaustive jusqu'à 40.

Les suites $\lfloor k\varphi \rfloor$ et $\lfloor k\varphi^2 \rfloor$ partitionnent exactement les entiers strictement positifs, d'après le théorème de Beatty\index{théorème de Beatty}, puisque $1/\varphi + 1/\varphi^2 = 1$. De plus, $\varphi^2 = \varphi + 1$, donc $\lfloor k\varphi^2 \rfloor = \lfloor k\varphi \rfloor + k$, qui est la différence requise. La formule respecte donc la règle de construction.

\iqwrong{Recourir de nouveau à la somme de Nim. Avec les coups diagonaux, ce n'est pas une somme de tas indépendants, et $9 \oplus 15 = 6$ n'est pas nul, alors même que la position est perdante.}
\end{solution}

\begin{solution}{iq:iv1:games:19}
La partie vaut $-1/18$ pour le joueur~1, donc le second joueur a l'avantage. Le joueur~1 dispose d'une famille entière d'équilibres. Pour $\alpha \in [0, \tfrac13]$, il mise le valet avec la probabilité $\alpha$, checke toujours la dame et suit une mise avec la probabilité $\tfrac13 + \alpha$, et mise le roi avec la probabilité $3\alpha$, sinon il checke et suit. La stratégie d'équilibre du joueur~2 est la suivante~: avec le roi, il mise ou suit toujours~; avec la dame, il checke lorsque le joueur~1 a checké, et suit une mise avec la probabilité $\tfrac13$~; avec le valet, il mise avec la probabilité $\tfrac13$ après un check, et se couche face à une mise.

Beaucoup de choix sont dominés\index{poker de Kuhn}. Le roi ne se couche jamais, le valet ne suit jamais, et la dame est trop forte pour bluffer et trop faible pour miser pour la valeur. Le valet bluffe donc, et le roi piège parfois son adversaire.

Maintenant, on utilise l'indifférence\index{fréquence de bluff}. Le joueur~2 bluffe avec le valet avec la probabilité $b$ après un check. Le joueur~1 avec la dame fait face à cette mise, venant d'un roi (toujours) ou d'un valet ($b$). Suivre rapporte $\frac{-2 + 2b}{1 + b}$, et se coucher rapporte $-1$, donc $b = \tfrac13$.

La dame du joueur~2, qui suit une mise avec la probabilité $c$, doit maintenir le valet du joueur~1 indifférent. Checker avec le valet perd 1 contre les deux. Miser rapporte $\tfrac12\bigl((1 - c) - 2c\bigr) + \tfrac12(-2)$. Ainsi $c = \tfrac13$.

La dame du joueur~2 est indifférente à une mise lorsque les mises sont constituées de trois quarts de rois. C'est pourquoi le roi du joueur~1 mise avec la probabilité $3\alpha$ lorsque le valet bluffe avec la probabilité $\alpha$. Le taux de suivi de la dame du joueur~1, égal à $\tfrac13 + \alpha$, maintient le valet du joueur~2 indifférent.

Un programme linéaire sur le tableau $27 \times 64$ des stratégies pures donne $-1/18$. Chaque membre de la famille, testé en $\alpha = 0$, $0.1$, $0.2$ et $\tfrac13$, rapporte au moins cette valeur face aux 64 réponses, et la stratégie du joueur~2 limite à au plus cette valeur les 27 stratégies pures.

\iqwrong{Supposer que celui qui mise le premier a l'avantage. Agir en second vaut davantage ici, car on voit ce que fait l'autre joueur avant de décider. Supposer aussi qu'il existe un équilibre unique, alors que le joueur~1 en a toute une famille de même valeur.}
\iqwhy{Un avantage de $1/18$ d'une mise de départ vient entièrement de la position, les deux joueurs ne commettant aucune erreur. L'essentiel de ce qu'on peut gagner en pratique vient d'un écart de l'adversaire par rapport à ces fréquences.}
\end{solution}

\begin{solution}{iq:iv1:games:20}
X mise avec les mains inférieures à $0.1$ et supérieures à $0.7$, et checke entre les deux. Y suit avec les mains supérieures à $0.4$. La partie vaut $\$0.10$ pour X.

Supposons que Y suive au-dessus de $c$\index{poker continu}\index{fréquence de bluff}. Checker avec la main $x$ rapporte $2x - 1$. Miser rapporte $c$ grâce aux couchages de Y. Avec $x$ inférieur à $c$, chaque suivi perd 2, donc miser rapporte $c - 2(1 - c) = 3c - 2$, quelle que soit la valeur de $x$. Avec $x$ supérieur à $c$, miser rapporte $c + 2\bigl((x - c) - (1 - x)\bigr)$.

Trois conditions déterminent la solution~:
\begin{itemize}
\item Une main de bluff en $a$ est indifférente~: $2a - 1 = 3c - 2$.
\item Une main de valeur en $b$ est indifférente~: $2b - 1 = c + 2(2b - c - 1)$, d'où $b = (1 + c)/2$.
\item Y en $c$ est indifférent entre se coucher, ce qui perd 1, et suivre, ce qui rapporte 2 contre un bluff et fait perdre 2 contre une main de valeur~: $\frac{2a - 2(1 - b)}{a + 1 - b} = -1$, d'où $3a = 1 - b$.
\end{itemize}
La résolution donne $c = 0.4$, $a = 0.1$ et $b = 0.7$. Miser avec les toutes pires mains fonctionne, car elles perdent de toute façon en checkant, et ce sont celles dont la valeur à l'abattage est la plus basse.

Pour une mise de taille $B$, les mêmes étapes donnent $a = \frac{B}{(B + 1)(B + 4)}$, $c = \frac{B(B + 3)}{(B + 1)(B + 4)}$ et $b = \frac{B^2 + 4B + 2}{(B + 1)(B + 4)}$. La valeur de la partie égale $a$, et elle est maximale en $B = 2$, où elle vaut $1/9$. Une intégration numérique confirme la valeur pour $B = 0.5$, 1, 2 et 4, et une simulation des règles en $B = 1$ donne $0.10$. Pour la plage donnée, tout autre seuil de suivi donne à X au moins autant, ce qui confirme que $c$ est la meilleure réponse de Y.

\iqwrong{Bluffer avec des mains moyennes, parce qu'elles sont les plus proches de gagner. Le bluff a un gain constant, donc on bluffe avec les mains qui ont le moins à perdre en ne checkant pas, c'est-à-dire les pires. Les mains moyennes préfèrent checker et gagner à l'abattage lorsqu'elles le peuvent.}
\iqwhy{Le rapport entre les bluffs et les mises de valeur est ici de $0.1:0.3$. Un suivi coûte 1~\$ pour gagner un pot de 3~\$, et un bluff pour trois mises de valeur rend le suiveur indifférent. Toute taille de mise affichée doit se lire de la même façon, comme un mélange de force et de bluff.}
\end{solution}
